% Unit 115 terjemahan Bahasa Indonesia dari chapter8.tex baris 1483--1564.
% Sumber: Methods of Algebra, Volume 2, commit
% 9a5803ff77dd3257484cb177f851a73770a59dd3 (CC BY 4.0).
% stable-unit-id: o014.aljabr2.chapter8.bisimplicial-a
% Cakupan: Bagian 8.8 dari definisi objek multisimplisial hingga contoh tensor.

% segment-id: o014.li2.chapter8.bisimplicial-a.h001
\section{Objek bisimplisial}\label{sec:bisimplicial}
\hypertarget{o014.aljabr2.chapter8.bisimplicial-a}{ }
% segment-id: o014.li2.chapter8.bisimplicial-a.p001
Kita mulai dengan beberapa definisi umum. Untuk objek-objek
$[m_1],\ldots,[m_n]$ dari $\simpDelta$, mulai sekarang objek
$([m_1],\ldots,[m_n])$ dari $(\simpDelta)^n$ juga akan ditulis sebagai
$[m_1]\times\cdots\times[m_n]$ agar lebih ringkas.

% segment-id: o014.li2.chapter8.bisimplicial-a.d001
\begin{definition}
	\index{objek bisimplisial@objek bisimplisial (bisimplicial object)}
	Misalkan $\mathcal{C}$ kategori sebarang dan $n\in\Z_{\geq1}$. Sebuah
	funktor berbentuk $X:(\simpDelta^{\opp})^n\to\mathcal{C}$ (atau
	$\simpDelta^n\to\mathcal{C}$) disebut \emph{objek simplisial $n$-lipat}
	(atau \emph{objek kosimplisial $n$-lipat}) dalam $\mathcal{C}$;
	morfismenya ialah morfisme antarfunktor tersebut. Bila $n=2$,
	objek yang bersesuaian disebut objek bisimplisial (atau bikosimplisial).
	Nilai objek simplisial $n$-lipat (atau kosimplisial $n$-lipat) $X$ pada
	$[m_1]\times\cdots\times[m_n]$ ditulis sebagai
	$X_{m_1,\ldots,m_n}$ (atau $X^{m_1,\ldots,m_n}$).
\end{definition}

% segment-id: o014.li2.chapter8.bisimplicial-a.p002
Jadi, objek simplisial $n$-lipat ditentukan oleh keluarga objek
$X_{m_1,\ldots,m_n}$, dengan
$m_1,\ldots,m_n\in\Z_{\geq0}$, beserta morfisme mukanya
% segment-id: o014.li2.chapter8.bisimplicial-a.q001
\[ {}^k d_i: X_{m_1, \ldots, m_n} \to X_{\ldots, m_k - 1, \ldots}, \quad 1 \leq k \leq n, \quad 0 \leq i \leq m_k \]
dan morfisme degenerasinya
% segment-id: o014.li2.chapter8.bisimplicial-a.q002
\[ {}^k s_j: X_{m_1, \ldots, m_n} \to X_{\ldots, m_k + 1, \ldots}, \quad 1 \leq k \leq n, \quad 0 \leq j \leq m_k . \]
Morfisme-morfisme tersebut disyaratkan memenuhi
\eqref{eqn:simplicial-identity} untuk setiap $k$, dan morfisme yang
bersesuaian dengan nilai-nilai $k$ yang berbeda harus saling komutatif.
Ini tidak lain adalah versi multivariabel
Definisi~\ref{def:simplicial-obj}.

% segment-id: o014.li2.chapter8.bisimplicial-a.p003
Secara lebih umum, untuk setiap morfisme dalam $\simpDelta^n$
% segment-id: o014.li2.chapter8.bisimplicial-a.q003
\[ f:[m_1]\times\cdots\times[m_n]\to[m'_1]\times\cdots\times[m'_n], \]
terdapat tarikan balik yang bersesuaian
$f^*:X_{m'_1,\ldots,m'_n}\to X_{m_1,\ldots,m_n}$. Kasus objek
kosimplisial $n$-lipat bersifat dual.

% segment-id: o014.li2.chapter8.bisimplicial-a.eg001
\begin{example}\label{eg:standard-n-fold-sSet}
	\index[sym1]{Deltapq@$\Delta^{p_1, \ldots, p_n}$}
	Ambil $\mathcal{C}=\cate{Set}$. Seperti dalam
	Definisi~\ref{def:standard-simplicial-set}, kita dapat mendefinisikan
	himpunan simplisial $n$-lipat standar
	% segment-id: o014.li2.chapter8.bisimplicial-a.q004
	\[ \Delta^{p_1, \ldots, p_n} := \Hom_{\simpDelta^n}\left(\cdot, [p_1] \times \cdots \times [p_n] \right). \]
	Selanjutnya ambil $\mathcal{C}=\cate{Ab}$. Dengan meniru
	Definisi~\ref{def:free-simplicial-Ab}, untuk setiap himpunan simplisial
	$n$-lipat $K$ kita dapat mendefinisikan objek simplisial $n$-lipat
	$\Z K$ dalam $\cate{Ab}$. Memberikan morfisme
	$\Z\Delta^{p_1,\ldots,p_n}\to X$ sama artinya dengan memberikan suatu
	elemen dari $X_{p_1,\ldots,p_n}$.
\end{example}

% segment-id: o014.li2.chapter8.bisimplicial-a.p004
Jika $\mathcal{C}$ kategori kecil, kategori semua objek simplisial
$n$-lipat dinotasikan oleh $\cate{s}^n\mathcal{C}$. Jadi
$\cate{s}^1\mathcal{C}=\cate{s}\mathcal{C}$, sedangkan bila $n>1$,
% segment-id: o014.li2.chapter8.bisimplicial-a.q005
\[ \cate{s}^n\mathcal{C}=\cate{s}\bigl(\cate{s}^{n-1}\mathcal{C}\bigr)
=\cate{s}^{n-1}\bigl(\cate{s}\mathcal{C}\bigr), \]
dan seterusnya. Kasus objek kosimplisial $n$-lipat serupa.

% segment-id: o014.li2.chapter8.bisimplicial-a.p005
Objek semisimplisial dalam Catatan~\ref{rem:semisimplicial} jelas memiliki
versi $n$-lipat, yang diperoleh dengan menghapus morfisme degenerasi dan
syarat-syarat terkait dari definisi di atas.

% segment-id: o014.li2.chapter8.bisimplicial-a.d002
\begin{definition}\label{def:bisimplicial-diagonal}
	\index{funktor!diagonal (diagonal)}
	Dengan notasi di atas, \emph{funktor diagonal}
	$d:\cate{s}^n\mathcal{C}\to\cate{s}\mathcal{C}$ memetakan objek
	simplisial $n$-lipat $X$ ke objek simplisial
	% segment-id: o014.li2.chapter8.bisimplicial-a.q006
	\[ d(X)_n := X_{n, \ldots, n}, \]
	dengan morfisme muka dan degenerasi yang berturut-turut didefinisikan
	sebagai $d_i=\prod_k{}^kd_i$ dan $s_j=\prod_k{}^ks_j$. Pada tingkat
	morfisme, funktor ini didefinisikan dengan cara yang jelas. Secara
	ekuivalen, $d(X)$ adalah komposisi $X$ dengan pembenaman diagonal
	$\simpDelta^{\opp}\to(\simpDelta^{\opp})^n$. Kasus objek kosimplisial
	sepenuhnya dual.
\end{definition}

% segment-id: o014.li2.chapter8.bisimplicial-a.eg002
\begin{example}\label{eg:boxtimes-bisimplicial}
	Misalkan $(\mathcal{C},\otimes)$ kategori monoidal, misalnya
	$\cate{Set}$ dengan hasil kali $\times$. Untuk
	$X_1,\ldots,X_n\in\Obj(\cate{s}\mathcal{C})$, kita memperoleh sebuah
	objek dari $\cate{s}^n\mathcal{C}$ yang suku
	$(m_1,\ldots,m_n)$-nya ialah
	$X_{1,m_1}\otimes\cdots\otimes X_{n,m_n}$. Objek simplisial $n$-lipat
	ini dinotasikan oleh
	% segment-id: o014.li2.chapter8.bisimplicial-a.q007
	\[ X_1 \boxtimes \cdots \boxtimes X_n \in \Obj(\cate{s}^n \mathcal{C}). \]

	Objek tersebut jangan dikacaukan dengan
	$X_1\otimes\cdots\otimes X_n\in\Obj(\cate{s}\mathcal{C})$ dari
	Definisi~\ref{def:monoidal-sObj}; hubungan keduanya ialah
	% segment-id: o014.li2.chapter8.bisimplicial-a.q008
	\[ X_1 \otimes \cdots \otimes X_n = d(X_1 \boxtimes \cdots \boxtimes X_n). \]
	Contoh dasarnya diperoleh dari kategori monoidal
	$(\cate{Set},\times)$. Dalam hal ini
	$\Delta^{p_1}\boxtimes\cdots\boxtimes\Delta^{p_n}
	=\Delta^{p_1,\ldots,p_n}$, sedangkan
	$\Delta^{p_1}\otimes\cdots\otimes\Delta^{p_n}
	=\Delta^{p_1}\times\cdots\times\Delta^{p_n}$, yaitu himpunan
	simplisial yang diperoleh dengan mengambil hasil kali suku demi suku.
\end{example}

% segment-id: o014.li2.chapter8.bisimplicial-a.p006
Objek simplisial $n$-lipat berperan penting dalam teori homotopi dan
penerapannya. Bagian ini terutama membahas kategori abelian, dengan
fokus pada kasus khusus $n=2$. Langkah pertama ialah memperkenalkan suatu
varian dari Definisi~\ref{def:unnormalized-chain}.

% segment-id: o014.li2.chapter8.bisimplicial-a.p007
Ingatlah definisi bikompleks dalam Definisi~\ref{def:double-cplx}, atau lebih
tepatnya versi kompleks rantainya. Kategori bikompleks rantai dinotasikan
oleh $\cate{Ch}^2(\cdots)$, dan $\cate{Ch}^2_{\geq0}(\cdots)$ didefinisikan
serupa. Untuk objek bisemisimplisial, kita juga memperkenalkan notasi
% segment-id: o014.li2.chapter8.bisimplicial-a.q009
\[ \dhori_i := {}^1 d_i , \quad \dvert_i := {}^2 d_i, \]
untuk morfisme muka pada arah ``horizontal'' dan ``vertikal''. Untuk objek
bisimplisial, morfisme degenerasi dinotasikan secara serupa oleh
$\shori_j={}^1s_j$ dan $\svert_j={}^2s_j$.

% segment-id: o014.li2.chapter8.bisimplicial-a.d003
\begin{definition}
	\index{kompleks rantai tak ternormalisasi@kompleks rantai tak ternormalisasi}
	\index[sym1]{CX2@$\mathrm{C}^2 X$}
	\index{bikompleks@bikompleks}
	Misalkan $\mathcal{A}$ kategori aditif dan $X$ objek bisemisimplisial
	dalam $\mathcal{A}$. \emph{Bikompleks rantai tak ternormalisasi} (atau
	\emph{bikompleks rantai Moore}) yang bersesuaian terdiri atas
	$\bigl(X_{p,q}\bigr)_{p,q\in\Z_{\geq0}}$ beserta data
	% segment-id: o014.li2.chapter8.bisimplicial-a.q010
	\begin{align*}
		\phori_{p, q} & := \sum_{i=0}^p (-1)^i \dhori_i: X_{p, q} \to X_{p-1, q}, \\
		\pvert_{p, q} & := \sum_{i=0}^q (-1)^i \dvert_i: X_{p, q} \to X_{p, q-1}.
	\end{align*}
	Menurut konvensi, $X_{p,q}:=0$ bila
	$(p,q)\notin\Z_{\geq0}^2$. Bikompleks rantai yang diperoleh dinotasikan
	$\mathrm{C}^2X\in\Obj(\cate{Ch}^2(\mathcal{A}))$.
\end{definition}

% segment-id: o014.li2.chapter8.bisimplicial-a.p008
Untuk menunjukkan bahwa ini benar-benar sebuah bikompleks rantai, kita
harus memverifikasi
$\phori_{p-1,q}\phori_{p,q}=0$,
$\pvert_{p,q-1}\pvert_{p,q}=0$, dan
$\pvert_{p-1,q}\phori_{p,q}=\phori_{p,q-1}\pvert_{p,q}$. Dua identitas
pertama sama persis dengan Definisi~\ref{def:unnormalized-chain}; identitas
terakhir langsung mengikuti dari komutativitas $\dhori_i$ dan $\dvert_j$.
Dengan demikian diperoleh funktor aditif
% segment-id: o014.li2.chapter8.bisimplicial-a.q011
\[ \mathrm{C}^2: \cate{s}^2 \mathcal{A} \to \cate{Ch}^2(\mathcal{A}). \]

% segment-id: o014.li2.chapter8.bisimplicial-a.p009
Karena $(p,q)\notin\Z_{\geq0}^2$ mengakibatkan $X_{p,q}=0$, sekarang kita
dapat mendefinisikan kompleks rantai totalnya secara alami sebagai
% segment-id: o014.li2.chapter8.bisimplicial-a.q012
\[ (\tot X)_n := \bigoplus_{p+q = n} X_{p, q}, \quad
\begin{tikzcd}[column sep=large]
	(\tot X)_n \arrow[r, "{\partial_n}"] & (\tot X)_{n-1} \\
	X_{p,q} \arrow[hookrightarrow, u] \arrow[r, "{(\phori_{p, q}, \; (-1)^p \pvert_{p, q})}"' inner sep=0.8em] & X_{p-1, q} \oplus X_{p, q-1} . \arrow[hookrightarrow, u]
\end{tikzcd}\]
Argumen yang sama dengan Definisi~\ref{def:total-cplx} menunjukkan bahwa
$\tot X$ benar-benar kompleks rantai. Jadi komposisi funktor
$\tot\circ\mathrm{C}^2:\cate{s}^2\mathcal{A}\to
\cate{Ch}_{\geq0}(\mathcal{A})$ terdefinisi.
\index{kompleks total@kompleks total}

% segment-id: o014.li2.chapter8.bisimplicial-a.p010
Konstruksi yang sama dapat diterapkan pada objek simplisial $n$-lipat di
atas $\mathcal{A}$ untuk memperoleh funktor
$\mathrm{C}^n:\cate{s}^n\mathcal{A}\to\cate{Ch}^n(\mathcal{A})$ beserta
kompleks total rantai yang bersesuaian; lihat
Catatan~\ref{rem:k-fold-cplx}.

% segment-id: o014.li2.chapter8.bisimplicial-a.eg003
\begin{example}\label{eg:tensor-bisimplicial}
	Misalkan $\mathcal{A}$ memiliki struktur kategori monoidal sedemikian
	rupa sehingga bifunktor
	$\otimes:\mathcal{A}\times\mathcal{A}\to\mathcal{A}$ aditif dalam
	setiap variabel. Jika
	$C_1,C_2\in\Obj(\cate{Ch}_{\geq0}(\mathcal{A}))$, kita memperoleh
	secara alami bikompleks rantai $C_1\boxtimes C_2$ dengan suku
	$(p,q)$ berupa $C_{1,p}\otimes C_{2,q}$ dan dengan $\phori$ serta
	$\pvert$ masing-masing berasal dari $C_1$ dan $C_2$. Jika
	$A,B\in\Obj(\cate{s}\mathcal{A})$, terdapat hubungan sederhana antara
	bikompleks rantai $\mathrm{C}A\boxtimes\mathrm{C}B$ dan objek
	bisimplisial $A\boxtimes B$ dari Contoh~\ref{eg:boxtimes-bisimplicial}:
	% segment-id: o014.li2.chapter8.bisimplicial-a.q013
	\[ \mathrm{C}^2 (A \boxtimes B) = \mathrm{C}A \boxtimes \mathrm{C}B. \]
	Selain itu, $\cate{Ch}_{\geq0}(\mathcal{A})$ menjadi kategori
	monoidal dengan
	% segment-id: o014.li2.chapter8.bisimplicial-a.q014
	\[ C_1 \otimes C_2 := \tot(C_1 \boxtimes C_2), \]
	dan unitnya jelas $\munit$ yang ditempatkan pada suku berderajat nol.
\end{example}
